\subsection{Affordance Successor Measures for \texttt{cube-single}}

We consider the state-based \texttt{cube-single} manipulation dataset.  Let the
offline dataset contain transitions
\[
\mathcal{D}
=
\{(s_t,a_t,s_{t+1},m_t)\}_{t=1}^N,
\]
where \(s_t\in\mathcal{S}\) is the state observation, \(a_t\in\mathcal{A}\) is the
continuous action, \(s_{t+1}\) is the next observation, and \(m_t\in\{0,1\}\)
is the bootstrap mask.  The usual FB successor measure would place mass on
future states \(s'\).  For manipulation, we instead define a successor measure
over \emph{affordance events}
\[
\zeta_t = \eta(s_t,a_t,s_{t+1}) \in \mathcal{Z}_{\mathcal{A}},
\]
where \(\eta\) extracts object-control information from the transition.

For \texttt{cube-single}, a useful event representation is
\[
\zeta_t
=
\Big[
p_{t+1}^{o},
\Delta p_t^{o},
h_{t+1}^{o},
\psi_{t+1}^{o},
d_{t+1}^{eo},
c_{t+1},
g_{t+1}^{\mathrm{open}},
a_t
\Big],
\]
where \(p_{t+1}^{o}\) is the cube position, \(\Delta p_t^{o}=p_{t+1}^{o}-p_t^{o}\)
is the cube displacement, \(h_{t+1}^{o}\) is object height, \(\psi_{t+1}^{o}\)
is object yaw, \(d_{t+1}^{eo}=\|p_{t+1}^{e}-p_{t+1}^{o}\|\) is end-effector/object
distance, \(c_{t+1}\) is a contact or near-contact indicator, \(g_{t+1}^{\mathrm{open}}\)
is gripper opening, and \(a_t\) is optionally included to support action-cost
rewards.  If the first implementation should remain minimal, we can start with
\[
\zeta_t
=
\Big[
p_{t+1}^{o},
\Delta p_t^{o},
h_{t+1}^{o},
d_{t+1}^{eo},
a_t
\Big].
\]

The affordance successor measure of a reward-conditioned policy \(\pi_z\) is
\[
\Psi_{\mathcal{A}}^{\pi_z}(s,a,d\zeta)
=
\mathbb{E}^{\pi_z}
\left[
\sum_{k=0}^{\infty}
\gamma^k
\delta_{\zeta_{t+k}}(d\zeta)
\;\middle|\;
s_t=s,\ a_t=a
\right].
\]
Unlike the standard successor measure
\[
M^{\pi_z}(s,a,ds')
=
\mathbb{E}^{\pi_z}
\left[
\sum_{k=0}^{\infty}
\gamma^k
\delta_{s_{t+k+1}}(ds')
\;\middle|\;
s_t=s,\ a_t=a
\right],
\]
this object does not primarily model future state visitation.  It models future
object-control events.  Thus, it asks which manipulation affordances become
available after taking action \(a\) in state \(s\).

We approximate the measure by a low-rank factorization
\[
\Psi_{\mathcal{A}}^{\pi_z}(s,a,d\zeta)
\approx
F_\theta(s,a,z)^\top B_\theta(\zeta)\,\nu(d\zeta),
\]
where \(\nu\) is the empirical distribution of affordance events in the dataset,
\(B_\theta:\mathcal{Z}_{\mathcal{A}}\to\mathbb{R}^d\) is the affordance-event basis,
and \(F_\theta:\mathcal{S}\times\mathcal{A}\times\mathbb{R}^d\to\mathbb{R}^d\)
is the forward successor feature.

For a reward over affordance events,
\[
r:\mathcal{Z}_{\mathcal{A}}\to\mathbb{R},
\]
the corresponding task vector is
\[
z_r
=
\mathbb{E}_{\zeta\sim\nu}
\left[
r(\zeta)B_\theta(\zeta)
\right].
\]
The induced action-value function is
\[
Q_r^{\pi_{z_r}}(s,a)
=
\int r(\zeta)\Psi_{\mathcal{A}}^{\pi_{z_r}}(s,a,d\zeta)
\approx
F_\theta(s,a,z_r)^\top z_r.
\]
This gives the same zero-shot readout interface as FB, but the reward is now a
function of manipulation-relevant affordance events rather than raw future states.

\paragraph{Pure affordance-FB objective.}
The direct FB analogue replaces future-state reconstruction with affordance-event
reconstruction.  The Bellman equation is
\[
\Psi_{\mathcal{A}}^{\pi_z}(s,a,d\zeta)
=
\delta_{\eta(s,a,s^+)}(d\zeta)
+
\gamma m_t
\Psi_{\mathcal{A}}^{\pi_z}
(s^+,\pi_z(s^+,z),d\zeta),
\]
where \(s^+=s_{t+1}\).  With the low-rank model, a sampled objective is
\[
\mathcal{L}_{\mathrm{AFB}}
=
\mathbb{E}_{(s,a,s^+)\sim\mathcal{D},\ z,\ \bar\zeta\sim\nu}
\left[
\left(
F_\theta(s,a,z)^\top B_\theta(\bar\zeta)
-
\gamma m_t
\bar F(s^+,\pi_z(s^+,z),z)^\top \bar B(\bar\zeta)
\right)^2
\right]
-
2\,
\mathbb{E}_{(s,a,s^+)\sim\mathcal{D},\ z}
\left[
F_\theta(s,a,z)^\top B_\theta(\zeta^+)
\right],
\]
where
\[
\zeta^+ = \eta(s,a,s^+).
\]
As in FB, we add an orthonormality constraint on the basis:
\[
\mathcal{L}_{\mathrm{ortho}}
=
\left\|
\mathbb{E}_{\zeta\sim\nu}
\left[
B_\theta(\zeta)B_\theta(\zeta)^\top
\right]
-
I
\right\|_F^2.
\]
The full measure objective is
\[
\mathcal{L}_{\mathrm{measure}}
=
\mathcal{L}_{\mathrm{AFB}}
+
\lambda_{\mathrm{ortho}}\mathcal{L}_{\mathrm{ortho}}.
\]

\paragraph{Value-coupled affordance successor objective.}
The pure measure objective still evaluates a policy family.  To preserve the main
lesson from bilinear GCIQL, we instead train the affordance successor features
through an IQL-style value-coupled objective.  For a sampled task vector \(z\),
define the projected reward
\[
r_z(\zeta)
=
B_\theta(\zeta)^\top z.
\]
The corresponding critic is
\[
Q_\theta(s,a,z)
=
F_\theta(s,a,z)^\top z.
\]
Given a transition \((s,a,s^+)\), compute
\[
\zeta^+ = \eta(s,a,s^+),
\qquad
r_z^+ = B_\theta(\zeta^+)^\top z.
\]
The value function is trained by an expectile regression over dataset actions:
\[
\mathcal{L}_V
=
\mathbb{E}_{(s,a)\sim\mathcal{D},\ z}
\left[
L_2^\tau
\left(
\bar Q(s,a,z)-V_\psi(s,z)
\right)
\right],
\]
where
\[
L_2^\tau(u)
=
|\tau-\mathbbm{1}[u<0]|u^2,
\qquad
\tau>\frac{1}{2}.
\]
The critic target is
\[
y_z
=
r_z^+
+
\gamma m_t V_{\bar\psi}(s^+,z),
\]
and the critic loss is
\[
\mathcal{L}_Q
=
\mathbb{E}_{(s,a,s^+)\sim\mathcal{D},\ z}
\left[
\left(
F_\theta(s,a,z)^\top z
-
y_z
\right)^2
\right].
\]
The representation should still be regularized so that \(B_\theta\) remains a useful
reward basis:
\[
\mathcal{L}_{B}
=
\lambda_{\mathrm{ortho}}
\left\|
\mathbb{E}_{\zeta\sim\nu}
\left[
B_\theta(\zeta)B_\theta(\zeta)^\top
\right]
-
I
\right\|_F^2
+
\lambda_{\mathrm{scale}}
\mathbb{E}_{\zeta\sim\nu}
\left[
\|B_\theta(\zeta)\|_2^2
\right].
\]
The value-coupled affordance-SM objective is
\[
\mathcal{L}_{\mathrm{ASM}}
=
\mathcal{L}_Q
+
\lambda_V\mathcal{L}_V
+
\mathcal{L}_{B}.
\]
This objective trains the successor representation and the value readout jointly.
It is therefore closer to bilinear GCIQL than to pure FB, but it retains FB's
zero-shot reward interface through \(z_r=\mathbb{E}[r(\zeta)B(\zeta)]\).

\paragraph{Task-vector distribution.}
The policy must be trained on the same kind of task vectors that will be used at
test time.  We therefore sample \(z\) from a mixture distribution
\[
p(z)
=
\alpha_{\mathrm{goal}}p_{\mathrm{goal}}(z)
+
\alpha_{\mathrm{rand}}p_{\mathrm{rand}}(z)
+
\alpha_{\mathrm{hand}}p_{\mathrm{hand}}(z),
\qquad
\alpha_{\mathrm{goal}}+\alpha_{\mathrm{rand}}+\alpha_{\mathrm{hand}}=1.
\]
Goal-reaching vectors are obtained from goal kernels:
\[
r_g(\zeta)
=
K(p^o,g),
\qquad
z_g
=
\mathbb{E}_{\zeta\sim\nu}
\left[
K(p^o,g)B_\theta(\zeta)
\right].
\]
Random reward vectors are sampled as
\[
z\sim \mathcal{N}(0,I),
\qquad
z\leftarrow \frac{z}{\|z\|_2}.
\]
Hand-designed manipulation reward vectors are obtained by first defining
\(r_i(\zeta)\) and then computing
\[
z_i
=
\mathbb{E}_{\zeta\sim\nu}
\left[
r_i(\zeta)B_\theta(\zeta)
\right].
\]
Useful hand-designed rewards include
\[
r_{\mathrm{right}}(\zeta)=\Delta p_x^o,
\]
\[
r_{\mathrm{lift}}(\zeta)=p_z^o,
\]
\[
r_{\mathrm{contact}}(\zeta)=\mathbbm{1}[d^{eo}<\epsilon],
\]
\[
r_{\mathrm{stable}}(\zeta)
=
\mathbbm{1}[d^{eo}<\epsilon]
-
\lambda_{\mathrm{slip}}\|\Delta p^o\|_2^2,
\]
\[
r_{\mathrm{avoid}}(\zeta)
=
-\mathbbm{1}[p^o\in\mathcal{B}_{\mathrm{bad}}],
\]
and, if yaw is included,
\[
r_{\mathrm{yaw}}(\zeta)
=
\cos(\psi^o-\psi^\star).
\]

\paragraph{Actor.}
We train a \(z\)-conditioned actor
\[
\pi_\omega(a\mid s,z)
\]
to realize the in-support maximizer of the learned affordance value.  A flow actor is
preferred because the action distribution in manipulation is multimodal.  Abstractly,
the actor objective is
\[
\max_\omega
\mathbb{E}_{s\sim\mathcal{D},\ z\sim p(z),\ a\sim \pi_\omega(\cdot\mid s,z)}
\left[
Q_\theta(s,a,z)
\right]
-
\beta_{\mathrm{BC}}\mathcal{L}_{\mathrm{BC}}(\omega),
\]
where
\[
Q_\theta(s,a,z)=F_\theta(s,a,z)^\top z.
\]
For a DDPG+BC-style actor, this can be written as
\[
\mathcal{L}_{\pi}
=
-
\mathbb{E}_{s,z}
\left[
Q_\theta(s,\pi_\omega(s,z),z)
\right]
+
\beta_{\mathrm{BC}}
\mathbb{E}_{(s,a)\sim\mathcal{D},z}
\left[
\|\pi_\omega(s,z)-a\|_2^2
\right].
\]
For a flow actor, \(\mathcal{L}_{\mathrm{BC}}\) is replaced by the flow matching or
flow-BC objective, and the Q-term is used to bias generated in-support actions toward
high affordance value.

\paragraph{Zero-shot policy extraction.}
At test time, define a new manipulation reward
\[
r:\mathcal{Z}_{\mathcal{A}}\to\mathbb{R}.
\]
Compute
\[
z_r
=
\mathbb{E}_{\zeta\sim\nu}
\left[
r(\zeta)B_\theta(\zeta)
\right].
\]
Then use
\[
a_t\sim \pi_\omega(\cdot\mid s_t,z_r),
\]
or choose the best action among flow samples:
\[
a_t
=
\arg\max_{a_i\sim\pi_\omega(\cdot\mid s_t,z_r)}
F_\theta(s_t,a_i,z_r)^\top z_r.
\]
This recovers FB-style zero-shot control, but over affordance rewards.

\paragraph{Policies enabled beyond bilinear GCIQL.}
A bilinear GCIQL critic trained as
\[
Q(s,a,g)=f(s,a)^\top \phi(g)
\]
primarily gives policies for reaching a specified goal \(g\), such as placing the
cube at a target position.  The affordance successor measure instead gives
\[
Q(s,a,z_r)=F(s,a,z_r)^\top z_r,
\]
where \(z_r\) may correspond to any reward over the affordance event space.  This
enables zero-shot extraction of policies such as
\[
\text{push the cube right:}\qquad r(\zeta)=\Delta p_x^o,
\]
\[
\text{lift the cube:}\qquad r(\zeta)=p_z^o,
\]
\[
\text{maintain contact:}\qquad r(\zeta)=\mathbbm{1}[d^{eo}<\epsilon],
\]
\[
\text{avoid a forbidden region:}\qquad
r(\zeta)=-\mathbbm{1}[p^o\in\mathcal{B}_{\mathrm{bad}}],
\]
\[
\text{stabilize the object:}\qquad
r(\zeta)=-\|\Delta p^o\|_2^2,
\]
\[
\text{reach a target while preserving contact:}\qquad
r(\zeta)=K(p^o,g)+\lambda\mathbbm{1}[d^{eo}<\epsilon],
\]
and, if the event representation contains yaw,
\[
\text{align orientation:}\qquad
r(\zeta)=\cos(\psi^o-\psi^\star).
These are not all naturally expressible as a single goal input \(g\).  They are
\]
reward functionals over object-control events.  The additional capability of the
affordance successor measure is therefore not merely better goal reaching; it is a
broader reward interface for manipulation.
